\documentclass[11pt,a4paper]{article}

\usepackage{geometry}
\geometry{a4paper, margin=1in}
\usepackage{xcolor}
\usepackage{colortbl}
\usepackage{booktabs}
\usepackage{tabularx}
\usepackage{enumitem}
\usepackage{setspace}
\usepackage{fancyhdr}
\usepackage{titlesec}

% Minimal fonts for XeLaTeX
\usepackage{fontspec}
\setmainfont{Calibri}

% Colors based on AIGGPA theme from previous doc
\definecolor{navyblue}{RGB}{0, 51, 102}
\definecolor{gold}{RGB}{204, 153, 51}
\definecolor{sectionbg}{RGB}{240, 244, 248}

\pagestyle{fancy}
\fancyhf{}
\fancyhead[L]{\footnotesize\color{navyblue}\textbf{AIGGPA Bhopal} \quad\textbar\quad Internal Pitch Document \& Q\&A}
\fancyfoot[C]{\footnotesize\color{navyblue}\thepage}
\renewcommand{\headrulewidth}{0.4pt}
\renewcommand{\headrule}{\hbox to\headwidth{\color{navyblue}\leaders\hrule height \headrulewidth\hfill}}

\titleformat{\section}
  {\normalfont\Large\bfseries\color{navyblue}}
  {\thesection}{1em}{}
\titleformat{\subsection}
  {\normalfont\large\bfseries\color{navyblue}}
  {\thesubsection}{1em}{}

\setlength{\headheight}{22pt}

\begin{document}

\begin{center}
    \vspace*{1cm}
    \Huge\color{navyblue}\textbf{RESEARCH PITCH \& STRATEGIC Q\&A} \\[0.5em]
    \large\color{gray}Assessment of the Use of Digital Tools by Government Employees \\[0.2cm]
    \normalsize\color{gold}\textbf{Prepared by:} Aryan Kori
    \vspace{1cm}
\end{center}

\section*{Part 1: The Oral Pitch (Speaking Notes)}
\textit{Use this as your spoken script when presenting the proposal to the management panel.}

\vspace{0.3cm}
\textbf{The Hook:} \\
"Good morning. Over the last few years, the government has invested heavily in digital infrastructure. But having the software is not the same as using it efficiently. My research proposal addresses the critical gap between top-down IT deployment and grassroots reality."

\vspace{0.2cm}
\textbf{The Problem:} \\
"Right now, employees in key departments like Revenue and Health are struggling. While central data shows high digital adoption, the reality at the local block level is that slow internet, difficult software, and lack of training often force government employees to revert to manual paper processing. We don't fully understand \textit{why} these block-level delays are happening."

\vspace{0.2cm}
\textbf{The Solution:} \\
"My study will use a targeted, qualitative approach. I will personally travel to district and block-level offices and speak directly with employees across all levels. Instead of relying on abstract statistics, I will collect firsthand data on exactly what is breaking down in their daily workflows."

\vspace{0.2cm}
\textbf{The Impact:} \\
"By mapping out these specific frustrations, this Institute will be able to tell the government exactly how to fix their training programs and streamline their IT helpdesks. This means less money wasted on unused software, highly confident employees, and ultimately - faster service for the citizens."

\vspace{0.5cm}
\hrule
\vspace{0.5cm}

\section*{Part 2: Anticipated Panel Questions (Defensive Q\&A)}
\textit{If the panel challenges the need for your research, use these answers to defend the proposal.}

\vspace{0.3cm}
\textbf{Question 1: "We already have IT data that shows high software adoption rates. Why do we need you to go into the field?"} \\
\textbf{Your Answer:} "Digital log-ins do not prove efficiency. An employee might log into a system because they are forced to, but it might be taking them twice as long to do the work because the system is confusing. My research will uncover the hidden efficiency losses that raw platform data cannot show."

\vspace{0.3cm}
\textbf{Question 2: "Why use qualitative interviews instead of sending out a mass digital survey to 5,000 employees?"} \\
\textbf{Your Answer:} "Frontline government employees are highly unlikely to document their genuine frustrations with state software on an official digital survey out of fear of getting in trouble. Physical, conversational interviews are the only way to build trust and capture authentic, unfiltered operational bottlenecks."

\vspace{0.3cm}
\textbf{Question 3: "Is 4 weeks enough time to gather meaningful data from multiple departments?"} \\
\textbf{Your Answer:} "Yes, because the methodology is highly structured. By using the 'Fieldwork Sample Collection Framework' included in my proposal, I have a targeted quota system. I am not trying to interview everyone - I am isolating specific employee cadres systematically to generate a precise heatmap of the issues."

\newpage

\section*{Part 3: Questions You Should Ask AIGGPA (Offensive Q\&A)}
\textit{Ask these questions \textbf{before} moving forward to show leadership that you are highly organized and proactive.}

\vspace{0.3cm}
\begin{itemize}[itemsep=15pt]
    \item \textbf{Regarding Access:} "Will the Institute provide me with an official authorization letter? I will need this to ensure that district officers are comfortable granting me access and taking time away from their work to be interviewed."
    
    \item \textbf{Regarding Scope:} "Are there specific under-performing districts or blocks that you want me to prioritize for my field visits? Or should I select the geographic sample randomly?"

    \item \textbf{Regarding The Final Product:} "Who is the primary stakeholder for the final report? Is it meant for internal AIGGPA policy-making, or will it be forwarded directly to the state IT department? Knowing this will help me tailor the vocabulary of my findings."
    
    \item \textbf{Regarding Budget \& Logistics:} "What is the standard procedure for claiming travel allowances for the fieldwork phase, and should I pre-clear the transportation modes for the rural block visits?"
\end{itemize}

\end{document}
